\section*{अध्याय \ref{ch:b2:frontier-orbitals} --- अग्रांत कक्षक और अभिक्रियाशीलता}

\begin{solution}{exo:b2:frontier-orbitals:1}
एथीन: \omterm{def:b2:frontier-orbitals:frontier}{HOMO} $\alpha + \beta$, \omterm{def:b2:frontier-orbitals:frontier}{LUMO} $\alpha - \beta$। ऐलिल ऋणायन (चार इलेक्ट्रॉन): \omterm{def:b2:frontier-orbitals:frontier}{HOMO} $\alpha$ (अनाबंधी
स्तर), \omterm{def:b2:frontier-orbitals:frontier}{LUMO} $\alpha - 1.414\beta$। ब्यूटाडाईन: \omterm{def:b2:frontier-orbitals:frontier}{HOMO} $\alpha + 0.618\beta$, \omterm{def:b2:frontier-orbitals:frontier}{LUMO} $\alpha -
0.618\beta$।
\end{solution}

\begin{solution}{exo:b2:frontier-orbitals:2}
कठोर: \ce{OH-}, \ce{H+}, \ce{F-}। मृदु: \ce{I-}, \ce{RS-}, आयडोमेथेन का कार्बन।
\end{solution}

\begin{solution}{exo:b2:frontier-orbitals:3}
ब्यूटाडाईन (s-\textit{cis} संरूपण सुलभ) और साइक्लोपेंटाडाईन (s-\textit{cis} में बँधी) अभिक्रिया करते हैं;
पेंटा-1,4-डाईन संयुग्मित नहीं है और अभिक्रिया नहीं करती। (2\textit{Z},4\textit{Z}) \omterm{def:b2:frontier-orbitals:diels-alder}{डाईन}
s-\textit{cis} रूप तक केवल अपने दोनों भीतरी मेथिल समूहों को एक-दूसरे में धकेलकर पहुँच सकती है:
वह बहुत कम अभिक्रिया करती है।
\end{solution}

\begin{solution}{exo:b2:frontier-orbitals:4}
साइक्लोहेक्स-3-ईन-1-कार्बोनाइट्राइल: एक षट्सदस्यी वलय, \omterm{def:b2:frontier-orbitals:diels-alder}{डाईन} के भूतपूर्व C2
और C3 के बीच द्वि-आबंध, \ce{C#N} समूह \omterm{def:b2:frontier-orbitals:diels-alder}{डाईनरागी} से बने वलय के दोनों
कार्बनों में से एक पर।
\end{solution}

\begin{solution}{exo:b2:frontier-orbitals:5}
$\Delta\varepsilon = \qty{8}{eV}$: विक्षोभ से $2\beta^2/\Delta\varepsilon = \qty{0.25}{eV}$। यथातथ: निचला स्तर
$-6 - \sqrt{16 + 1} = \qty{-10.123}{eV}$, लाभ $2 \times 0.123 = \qty{0.246}{eV}$।
\end{solution}

\begin{solution}{exo:b2:frontier-orbitals:6}
आयडोमेथेन \omterm{def:b2:frontier-orbitals:hard-soft}{मृदु इलेक्ट्रॉनरागी} है: \omterm{def:b2:frontier-orbitals:control}{कक्षक नियंत्रण}, आक्रमण उस परमाणु से जिसका
\omterm{def:b2:frontier-orbitals:frontier}{HOMO} गुणांक बड़ा है, अर्थात् कार्बन। मुक्त कार्बधनायन \omterm{def:b2:frontier-orbitals:hard-soft}{कठोर इलेक्ट्रॉनरागी} है: \omterm{def:b2:frontier-orbitals:control}{आवेश नियंत्रण},
आक्रमण अंशतः नाइट्रोजन से, जिस पर ऋणात्मक आवेश अधिक है।
\end{solution}

\begin{solution}{exo:b2:frontier-orbitals:7}
\omterm{def:b2:frontier-orbitals:diels-alder}{डाईन} का C4 (सबसे बड़ा \omterm{def:b2:frontier-orbitals:frontier}{HOMO} गुणांक) प्रोपीनल के $\beta$ कार्बन (सबसे बड़ा
\omterm{def:b2:frontier-orbitals:frontier}{LUMO} गुणांक) से आबंधित होता है: 2-मेथॉक्सीसाइक्लोहेक्स-3-ईन-1-कार्बैल्डिहाइड, दोनों प्रतिस्थापी
पड़ोसी कार्बनों पर (``ऑर्थो'' उत्पाद)।
\end{solution}

\begin{solution}{exo:b2:frontier-orbitals:8}
बाइसाइक्लो[2.2.1]हेप्टीन (नॉरबोर्नीन) ढाँचा, जिसमें एस्टर समूह भूतपूर्व \omterm{def:b2:frontier-orbitals:diels-alder}{डाईनरागी} के
एक कार्बन पर है, नए \ce{C=C} आबंध की ओर, षट्सदस्यी वलय के नीचे
(endo), \ce{CH2} सेतु की विपरीत ओर इंगित करता हुआ।
\end{solution}

\begin{solution}{exo:b2:frontier-orbitals:9}
दोनों कार्बोनिल समूह \omterm{def:b2:frontier-orbitals:diels-alder}{डाईनरागी} के \omterm{def:b2:frontier-orbitals:frontier}{LUMO} को नीचे लाते हैं: साइक्लोपेंटाडाईन के \omterm{def:b2:frontier-orbitals:frontier}{HOMO} से उसका
अंतराल एथीन की तुलना में बहुत छोटा है, संक्रमण अवस्था का स्थायीकरण $\beta^2/\Delta\varepsilon$
बड़ा है, और अवरोध नीचा।
\end{solution}

\begin{solution}{exo:b2:frontier-orbitals:10}
अभिक्रिया त्रिविमविशिष्ट है: मैलिएट वह योगज देता है जिसमें दोनों एस्टर समूह cis हैं
(मुख्य उत्पाद के रूप में endo,endo; यह अकिरैल है, एक मेसो यौगिक); फ़्यूमेरेट
trans योगज देता है, एक एस्टर endo और एक exo, जो प्रतिबिंबरूपों के युग्म (रेसिमिक) के रूप में बनता है।
\end{solution}

\begin{solution}{exo:b2:frontier-orbitals:11}
दो $\pi$ आबंध दो अधिक प्रबल $\sigma$ आबंध बन जाते हैं: $\Delta_r H^\circ < 0$। दो अणु एक बन जाते हैं:
$\Delta_r S^\circ < 0$। $\Delta_r G^\circ = \Delta_r H^\circ - T\Delta_r S^\circ$ $T$ के साथ बढ़ता है और
$T_i = \Delta_r H^\circ/\Delta_r S^\circ$ से ऊपर धनात्मक हो जाता है: उच्च ताप पर योगज वापस टूट जाता है
(प्रतीप डील्स--ऐल्डर)।
\end{solution}

\begin{solution}{exo:b2:frontier-orbitals:12}
एकाकी युग्मों के दो भरे कक्षक पास-पास चार-इलेक्ट्रॉन अन्योन्यक्रिया बनाते हैं, जो
प्रतिकर्षित करती है। अणु \ce{N-N} आबंध के परितः तब तक घूमता है जब तक एकाकी युग्म लगभग
लंबवत न हो जाएँ (एक गॉश संरूपण), जहाँ उनका अतिव्यापन, और प्रतिकर्षण,
छोटा होता है।
\end{solution}

\begin{solution}{pb:b2:frontier-orbitals:1}
\textbf{1.} दो संयुग्मित द्वि-आबंधों और एक \ce{CH2} वाला पंचसदस्यी वलय: वलय
\omterm{def:b2:frontier-orbitals:diels-alder}{डाईन} इकाई को s-\textit{cis} संरूपण में बाँधे रखता है।
\textbf{2.} \omterm{def:b2:frontier-orbitals:diels-alder}{डाईनरागी}, अपने एक द्वि-आबंध के द्वारा।
\textbf{3.} साइक्लोपेंटीन वलय से संगलित नॉरबोर्नीन ढाँचा; दोनों नए $\sigma$ आबंध
\omterm{def:b2:frontier-orbitals:diels-alder}{डाईन} के सिरों को \omterm{def:b2:frontier-orbitals:diels-alder}{डाईनरागी} के द्वि-आबंध के दोनों कार्बनों से जोड़ते हैं।
\textbf{4.} \omterm{def:b2:frontier-orbitals:endo}{endo योगज} (\omterm{def:b2:frontier-orbitals:endo}{endo नियम}): संक्रमण अवस्था में \omterm{def:b2:frontier-orbitals:diels-alder}{डाईनरागी} का शेष वलय
\omterm{def:b2:frontier-orbitals:diels-alder}{डाईन} के नीचे रहता है।
\textbf{5.} हाँ: यह समन्वित है, दोनों आबंध हर भागीदार के एक ही फलक पर बनते हैं, और
\omterm{def:b2:frontier-orbitals:diels-alder}{डाईनरागी} की ज्यामिति बनी रहती है।
\textbf{6.} साइक्लोपेंटाडाईन एक अच्छी \omterm{def:b2:frontier-orbitals:diels-alder}{डाईन} (s-\textit{cis} में बँधी) भी है और \omterm{def:b2:frontier-orbitals:diels-alder}{डाईनरागी} भी, और
अभिक्रिया समन्वित है, जिसमें स्थायी करने के लिए कोई आवेशित मध्यवर्ती नहीं।
\textbf{7.} ऋणात्मक: दो $\pi$ आबंधों का स्थान दो $\sigma$ आबंध लेते हैं।
\textbf{8.} ऋणात्मक: दो अणु एक देते हैं।
\textbf{9.} $\Delta_r G^\circ = \Delta_r H^\circ - T\Delta_r S^\circ$, दोनों पद ऋणात्मक: निम्न
$T$ पर ऋणात्मक, $T_i = \Delta_r H^\circ/\Delta_r S^\circ$ पर शून्य, उससे ऊपर धनात्मक।
\textbf{10.} कमरे के ताप पर $T < T_i$: द्वितयन अनुकूल है, यद्यपि धीमा। द्वितय के
क्वथनांक के निकट साम्य एकलक के अनुकूल होता है, और वह शीघ्र
प्राप्त हो जाता है।
\textbf{11.} एकलक गर्म मिश्रण से वाष्प के रूप में निकल जाता है: किसी उत्पाद को हटाते रहने से
साम्य एकलक की ओर खिसकता रहता है (ले शातेलिए)।
\textbf{12.} कमरे के ताप पर वह फिर द्वितय बनाता है; ठंडक उस अभिक्रिया को धीमा करती है।
\textbf{13.} $0.62 - (-0.62) = 1.24|\beta|$।
\textbf{14.} $0.62 - (-0.35) = 0.97|\beta|$।
\textbf{15.} मैलेइक ऐनहाइड्राइड के साथ साइक्लोपेंटाडाईन: छोटा अंतराल संक्रमण अवस्था का
बड़ा स्थायीकरण देता है, अर्थात् तेज़ अभिक्रिया।
\textbf{16.} वे \ce{C=C} आबंध से संयुग्मित और इलेक्ट्रॉन-आकर्षी हैं: प्रोपीनल की तरह,
वे $\pi^*$ स्तर को नीचे खींचते हैं।
\textbf{17.} उसका \omterm{def:b2:frontier-orbitals:frontier}{LUMO}, जिसकी कार्बोनिल कार्बनों पर पालियाँ \omterm{def:b2:frontier-orbitals:diels-alder}{डाईन} के \omterm{def:b2:frontier-orbitals:frontier}{HOMO} के
केंद्रीय कार्बनों के सामने होती हैं।
\textbf{18.} नॉरबोर्नीन ढाँचा; ऐनहाइड्राइड वलय नए द्वि-आबंध के नीचे,
\ce{CH2} सेतु की विपरीत ओर।
\textbf{19.} दोनों सेतुशीर्ष कार्बन और ऐनहाइड्राइड धारण करने वाले दोनों कार्बन: चार
त्रिविमजनक केंद्र।
\textbf{20.} नहीं: एक दर्पण तल \ce{CH2} सेतु से होकर और अणु के दोनों अर्धांशों के बीच से
गुज़रता है; यह एक मेसो यौगिक है।
\textbf{21.} \omterm{def:b2:frontier-orbitals:diels-alder}{डाईनरागी} पर वे cis थे, और समन्वित योग उन्हें cis ही रखता है।
\textbf{22.} वही ढाँचा, ऐनहाइड्राइड वलय \ce{CH2} सेतु की ओर।
\textbf{23.} नहीं: endo को exo में बदलने का अर्थ होगा \ce{C-C} आबंधों को तोड़ना और फिर बनाना,
प्रतीप अभिक्रिया के द्वारा, जिसके लिए गर्म करना पड़ता है।
\textbf{24.} जल ऐनहाइड्राइड को खोल देता है: cis डाइकार्बोक्सिलिक अम्ल, दोनों \ce{COOH} समूह endo।
\textbf{25.} \omterm{def:b2:frontier-orbitals:endo}{endo योगज} के $\boldsymbol{4}$ त्रिविमजनक केंद्र हैं।
\end{solution}
